\documentclass[12pt,a4paper]{article}

\usepackage[utf8]{inputenc}
\usepackage[T2A]{fontenc}
\usepackage[russian,english]{babel}
\usepackage{geometry}
\usepackage{array}
\usepackage{longtable}
\usepackage{enumitem}
\usepackage{hyperref}
\usepackage{xcolor}

\geometry{margin=2cm}
\hypersetup{
    colorlinks=true,
    linkcolor=blue,
    urlcolor=blue
}

\setlist[itemize]{topsep=4pt,itemsep=2pt}
\setlist[enumerate]{topsep=4pt,itemsep=2pt}

\begin{document}

\selectlanguage{russian}

\begin{center}
    {\LARGE \textbf{Sprint 1}}\\[6pt]
    {\large План работ на неделю: Planning, Setup \& MVP Development}
\end{center}

\section{Цель спринта}

Цель Sprint 1 --- подготовить проектную базу, закрыть требования Week 2 и начать реализацию MVP для системы AI-интеграции в GitFlame.

В рамках спринта команда должна:

\begin{itemize}
    \item описать проект, выбранный трек, проблему и MVP scope;
    \item подготовить архитектуру решения;
    \item выбрать open-source модели для двух ML-сценариев;
    \item поднять backend, frontend, базу данных и Docker-инфраструктуру;
    \item реализовать первый end-to-end MVP flow;
    \item подготовить документацию, wireframes, backlog и материалы для отчёта.
\end{itemize}

\section{Ключевая идея проекта}

Проект расширяет GitFlame AI-функциональностью для работы с репозиториями. При этом команда не имеет прямого доступа к внутреннему коду GitFlame, поэтому в Sprint 1 мы проектируем не прямое изменение GitFlame, а внешний интеграционный сервис и контракты взаимодействия.

GitFlame в нашем MVP рассматривается как внешний продукт, который:

\begin{itemize}
    \item может отправлять события в нашу систему через API или webhook;
    \item может передавать данные issue, repository metadata и содержимое \texttt{.yml};
    \item может получать от нашей системы общие ответы: plan, status, recommendations, branch/PR request payload;
    \item может отображать результат в своём интерфейсе на основе API-ответов нашей системы.
\end{itemize}

Система должна поддерживать два основных сценария:

\begin{enumerate}
    \item \textbf{Автогенерация кода по issue.}
    Пользователь создаёт issue, система проверяет наличие конфигурационного файла \texttt{.yml}, отправляет данные в ML-сервис, получает план реализации в формате \texttt{.md}, ожидает реакцию пользователя и при подтверждении создаёт ветку, файлы и pull request.

    \item \textbf{Рекомендательная система.}
    При наличии \texttt{.yml} система анализирует репозиторий, формирует краткое summary и список рекомендаций, сохраняет результат в базе данных и отображает его в отдельном виджете на вкладке Code.
\end{enumerate}

\section{Архитектурное решение}

В Sprint 1 необходимо заложить архитектуру, в которой сценарии автогенерации и рекомендаций могут использовать разные open-source модели.

\begin{itemize}
    \item \texttt{code\_generation\_model} --- модель для генерации плана и кода по issue.
    \item \texttt{recommendation\_model} --- модель для анализа репозитория и генерации рекомендаций.
    \item Выбор моделей является внутренней настройкой нашей системы и не задаётся пользователем в \texttt{.yml}.
    \item Формат плана всегда \texttt{Markdown}; пользователь не выбирает другой формат.
\end{itemize}

Общий поток автогенерации:

\begin{center}
\texttt{GitFlame UI -> Backend -> ML Service N -> Backend -> Git Workflow -> Pull Request}
\end{center}

Поток рекомендаций:

\begin{center}
\texttt{Repository -> .yml -> ML Analysis -> Database -> Recommendations Widget}
\end{center}

\section{Интеграционные ограничения}

Так как у команды нет доступа к внутреннему коду GitFlame, в Sprint 1 необходимо явно зафиксировать границы системы:

\begin{itemize}
    \item мы не реализуем реальные изменения внутри GitFlame;
    \item мы не создаём настоящий production widget внутри GitFlame UI;
    \item мы проектируем API, JSON-контракты и mock UI, который демонстрирует ожидаемое поведение;
    \item мы ожидаем, что GitFlame будет передавать нашей системе \texttt{.yml}, данные issue и metadata репозитория;
    \item мы возвращаем GitFlame структурированные ответы, которые он сможет использовать для комментариев, статусов, рекомендаций и PR workflow;
    \item Git branch/PR workflow в Sprint 1 реализуется как mock или contract-level service, потому что реальный GitFlame API может быть недоступен.
\end{itemize}

\section{Предварительная структура \texttt{.yml}}

Структуру конфигурационного файла команда должна продумать самостоятельно. В Sprint 1 необходимо подготовить draft спецификации \texttt{.yml}, который будет использоваться backend и ML-сервисом.

Важно: \texttt{.yml} хранит только настройки анализа и поведения системы, доступные пользователю. Он не содержит выбор ML-моделей и не содержит выбор формата плана. Модели выбираются нашей системой, а план всегда возвращается в формате \texttt{Markdown}.

Пример базовой структуры:

\begin{verbatim}
version: 1

repository:
  default_branch: main
  target_branch_prefix: ai/

analysis:
  enabled: true
  include:
    - src/**
    - internal/**
  exclude:
    - node_modules/**
    - dist/**
    - build/**
    - .git/**

code_generation:
  enabled: true
  require_user_approval: true
  reviewer_policy: issue_author
  allowed_actions:
    approve_command: "/approve"
    correct_command: "/correct"
    reject_command: "/reject"

recommendations:
  enabled: true
  severity_threshold: low
  categories:
    - code_duplication
    - security
    - maintainability
    - performance
    - architecture

rag:
  max_files: 20
  max_file_size_kb: 120
  context_strategy: issue_relevant_files

storage:
  recommendation_ttl_days: 30
\end{verbatim}

В Sprint 1 эта структура считается draft-версией. После уточнения требований заказчика она может быть изменена.

\section{Распределение задач}

\subsection{Данил}

\textbf{Роль: Team Lead + ML}

\textbf{Основная зона ответственности:} архитектура, координация спринта, ML-сценарий автогенерации кода.

\textbf{Задачи:}

\begin{enumerate}
    \item Подготовить \texttt{rubric\_checklist.md}:
    \begin{itemize}
        \item выписать все требования Week 2;
        \item указать, какой артефакт закрывает каждое требование;
        \item проверить в конце недели, что все deliverables готовы.
    \end{itemize}

    \item Описать общую архитектуру MVP:
    \begin{itemize}
        \item схема автогенерации кода;
        \item схема рекомендательной системы;
        \item разделение моделей для автогенерации и рекомендаций.
        \item границы интеграции с GitFlame без доступа к его внутреннему коду.
    \end{itemize}

    \item Описать Industrial track justification:
    \begin{itemize}
        \item existing product: GitFlame;
        \item gap: отсутствие AI-автогенерации кода и AI-рекомендаций;
        \item contribution: \texttt{.yml}-конфиг, issue-to-plan workflow, approve/correct/reject, PR generation, recommendation system;
        \item impact: ускорение обработки issue и улучшение анализа репозитория.
    \end{itemize}

    \item Определить MVP scope:
    \begin{itemize}
        \item IN: draft \texttt{.yml} specification, issue-to-plan, approve/correct/reject, API contracts, mock branch/PR, recommendations summary/cards, demo UI;
        \item OUT: production-ready ML deployment, идеальная генерация кода, сложная auth-система, полноценная RAG-инфраструктура для больших репозиториев, реальные изменения во внутреннем коде GitFlame.
    \end{itemize}

    \item Продумать структуру \texttt{.yml}:
    \begin{itemize}
        \item описать поля для repository settings;
        \item описать поля для analysis include/exclude;
        \item описать поля для code generation;
        \item описать поля для recommendations;
        \item описать поля для RAG/context retrieval;
        \item явно зафиксировать, что пользователь не выбирает ML-модель;
        \item явно зафиксировать, что план всегда возвращается в Markdown;
        \item подготовить пример валидного \texttt{.yml}.
    \end{itemize}

    \item Создать Kanban backlog:
    \begin{itemize}
        \item завести задачи для всех участников;
        \item расставить приоритеты;
        \item добавить labels: \texttt{ML}, \texttt{Backend}, \texttt{Frontend}, \texttt{Docs}, \texttt{Infra}, \texttt{MVP}.
    \end{itemize}

    \item Выполнить ML-часть для автогенерации:
    \begin{itemize}
        \item выбрать open-source модель для автогенерации;
        \item основной кандидат: \texttt{Qwen2.5-Coder-32B-Instruct};
        \item fallback: \texttt{Qwen2.5-Coder-14B-Instruct} или quantized 32B;
        \item написать \texttt{autogen\_prompt.md};
        \item описать формат \texttt{plan.md};
        \item описать RAG/context retrieval;
        \item описать JSON request/response для ML-сервиса автогенерации.
    \end{itemize}

    \item Провести internal demo:
    \begin{itemize}
        \item показать текущий MVP flow;
        \item собрать список багов;
        \item записать next steps.
    \end{itemize}
\end{enumerate}

\textbf{Deliverables:}
\begin{itemize}
    \item \texttt{rubric\_checklist.md}
    \item \texttt{architecture.md}
    \item \texttt{mvp\_scope.md}
    \item \texttt{ai\_config\_spec.md}
    \item example \texttt{.yml}
    \item \texttt{autogen\_prompt.md}
    \item \texttt{rag\_strategy.md}
    \item model choice для автогенерации
    \item Kanban backlog
    \item notes from internal demo
\end{itemize}

\subsection{Карим}

\textbf{Роль: ML}

\textbf{Основная зона ответственности:} рекомендательная ML-система и выбор модели для анализа репозитория.

\textbf{Задачи:}

\begin{enumerate}
    \item Выбрать open-source модель для рекомендательной системы:
    \begin{itemize}
        \item основной кандидат: \texttt{Qwen2.5-Coder-14B-Instruct};
        \item fallback: \texttt{Qwen2.5-Coder-7B-Instruct};
        \item вариант для мощного сервера: \texttt{Qwen2.5-Coder-32B-Instruct};
        \item сравнить модели по качеству анализа кода, требованиям к железу, лицензии и скорости.
    \end{itemize}

    \item Подготовить \texttt{recommendation\_prompt.md}:
    \begin{itemize}
        \item prompt для анализа репозитория;
        \item учёт параметров из \texttt{.yml};
        \item возврат summary и списка рекомендаций.
    \end{itemize}

    \item Описать JSON schema для рекомендаций:
    \begin{itemize}
        \item \texttt{summary};
        \item \texttt{recommendations[]};
        \item \texttt{severity};
        \item \texttt{file};
        \item \texttt{line};
        \item \texttt{problem};
        \item \texttt{suggestion};
        \item \texttt{confidence}.
    \end{itemize}

    \item Подготовить mock ML service для рекомендаций:
    \begin{itemize}
        \item input: \texttt{.yml + repo context};
        \item output: \texttt{summary + recommendation cards};
        \item endpoint должен использовать финальный формат ответа.
    \end{itemize}

    \item Подготовить ML experimentation setup:
    \begin{itemize}
        \item папка, скрипт или notebook для тестов;
        \item несколько тестовых файлов;
        \item пример запуска анализа.
    \end{itemize}

    \item Сделать initial baseline:
    \begin{itemize}
        \item взять небольшой пример репозитория или файлов;
        \item прогнать рекомендательную модель или mock;
        \item показать первый результат;
        \item описать сравнение с ручным анализом.
    \end{itemize}
\end{enumerate}

\textbf{Deliverables:}
\begin{itemize}
    \item \texttt{recommendation\_prompt.md}
    \item \texttt{recommendation\_schema.json}
    \item model comparison для рекомендаций
    \item mock recommendation ML service
    \item experimentation setup
    \item initial baseline result
\end{itemize}

\subsection{Артур}

\textbf{Роль: Backend}

\textbf{Основная зона ответственности:} backend skeleton, API contract и issue workflow.

\textbf{Задачи:}

\begin{enumerate}
    \item Инициализировать backend boilerplate:
    \begin{itemize}
        \item создать структуру проекта;
        \item добавить базовую конфигурацию;
        \item реализовать \texttt{GET /health}.
    \end{itemize}

    \item Подготовить OpenAPI/Swagger:
    \begin{itemize}
        \item описать MVP endpoints;
        \item добавить request/response schemas;
        \item подготовить ссылку на API documentation.
    \end{itemize}

    \item Реализовать issue workflow endpoints для взаимодействия с внешним GitFlame:
    \begin{itemize}
        \item \texttt{POST /integrations/gitflame/issues/analyze};
        \item \texttt{GET /ai/issues/\{id\}/plan};
        \item \texttt{POST /ai/issues/\{id\}/approve};
        \item \texttt{POST /ai/issues/\{id\}/correct};
        \item \texttt{POST /ai/issues/\{id\}/reject}.
    \end{itemize}

    \item Реализовать интеграцию с ML service для автогенерации:
    \begin{itemize}
        \item принимать от GitFlame \texttt{.yml}, issue title/body и repo context;
        \item отправлять эти данные в ML-сервис;
        \item получать \texttt{plan.md};
        \item передавать результат в слой хранения;
        \item возвращать GitFlame структурированный response для комментария к issue.
    \end{itemize}

    \item Добавить базовую обработку ошибок:
    \begin{itemize}
        \item отсутствует \texttt{.yml};
        \item ML service недоступен;
        \item модель вернула некорректный ответ;
        \item session id или issue id не найден.
    \end{itemize}
\end{enumerate}

\textbf{Deliverables:}
\begin{itemize}
    \item backend skeleton
    \item runnable \texttt{GET /health}
    \item OpenAPI/Swagger
    \item issue workflow endpoints
    \item ML service client для автогенерации
    \item backend PR/Issue links
\end{itemize}

\subsection{Амир}

\textbf{Роль: Backend}

\textbf{Основная зона ответственности:} база данных и хранение данных MVP.

\textbf{Задачи:}

\begin{enumerate}
    \item Спроектировать initial DB schema:
    \begin{itemize}
        \item \texttt{repositories};
        \item \texttt{ai\_configs};
        \item \texttt{issue\_sessions};
        \item \texttt{generated\_plans};
        \item \texttt{user\_responses};
        \item \texttt{recommendations};
        \item \texttt{recommendation\_statuses}.
    \end{itemize}

    \item Реализовать модели или таблицы:
    \begin{itemize}
        \item подготовить migrations или schema-файл;
        \item добавить базовые entity/model классы.
    \end{itemize}

    \item Реализовать хранение issue workflow:
    \begin{itemize}
        \item сохранять issue session;
        \item сохранять generated plan;
        \item сохранять approve/correct/reject;
        \item хранить текущий статус сессии.
    \end{itemize}

    \item Реализовать хранение рекомендаций:
    \begin{itemize}
        \item сохранять summary;
        \item сохранять recommendation cards;
        \item обновлять статус рекомендации;
        \item удалять рекомендацию.
    \end{itemize}

    \item Проверить data flow:
    \begin{itemize}
        \item backend сохраняет и отдаёт план;
        \item backend сохраняет и отдаёт рекомендации.
    \end{itemize}
\end{enumerate}

\textbf{Deliverables:}
\begin{itemize}
    \item DB schema
    \item migrations/schema file
    \item storage layer
    \item saved issue sessions/plans
    \item saved recommendations
    \item basic DB verification
\end{itemize}

\subsection{Руслан}

\textbf{Роль: Backend / Infrastructure}

\textbf{Основная зона ответственности:} Docker, \texttt{.yml} config service и Git workflow.

\textbf{Задачи:}

\begin{enumerate}
    \item Подготовить Docker support:
    \begin{itemize}
        \item Dockerfile для backend;
        \item \texttt{docker-compose.yml} для backend и database;
        \item инструкция запуска.
    \end{itemize}

    \item Реализовать config service для \texttt{.yml}:
    \begin{itemize}
        \item принять \texttt{.yml} от внешнего GitFlame API/webhook;
        \item прочитать \texttt{.yml};
        \item провалидировать поля;
        \item применить branch, exclusions и rules.
    \end{itemize}

    \item Реализовать Git workflow service:
    \begin{itemize}
        \item подготовить contract для запроса на создание branch;
        \item подготовить contract для передачи сгенерированных файлов;
        \item подготовить contract для запроса на создание pull request;
        \item подготовить contract для назначения reviewer автора issue;
        \item сделать mock-ответы, так как реальный GitFlame API может быть недоступен.
    \end{itemize}

    \item Для MVP сделать mock Git workflow:
    \begin{itemize}
        \item возвращать fake branch name;
        \item возвращать fake PR URL;
        \item возвращать fake reviewer;
        \item сохранить интерфейс так, чтобы later заменить mock на реальный GitFlame/Git API.
    \end{itemize}
\end{enumerate}

\textbf{Deliverables:}
\begin{itemize}
    \item Dockerfile
    \item \texttt{docker-compose.yml}
    \item \texttt{.yml} config service
    \item Git workflow service interface
    \item mock branch/PR flow
    \item infra PR/Issue links
\end{itemize}

\subsection{Рома}

\textbf{Роль: Backend Support + Frontend + Documentation}

\textbf{Основная зона ответственности:} frontend skeleton, recommendation API support, документация и визуальные материалы.

\textbf{Задачи:}

\begin{enumerate}
    \item Подготовить report structure:
    \begin{itemize}
        \item team name;
        \item communication channels;
        \item members and roles;
        \item project statement;
        \item problem statement;
        \item track justification;
        \item roadmap;
        \item risks;
        \item MVP scope;
        \item user stories;
        \item acceptance criteria;
        \item implemented features;
        \item screenshots/GIFs;
        \item PR/Issue/API links;
        \item internal demo notes.
    \end{itemize}

    \item Написать user stories:
    \begin{itemize}
        \item пользователь хочет включить AI в репозитории;
        \item пользователь хочет настроить \texttt{.yml};
        \item пользователь хочет получить план по issue;
        \item пользователь хочет approve/correct/reject план;
        \item пользователь хочет видеть summary рекомендаций;
        \item пользователь хочет открыть подробный анализ;
        \item пользователь хочет закрыть или удалить рекомендацию.
    \end{itemize}

    \item Сделать wireframes/mockups:
    \begin{itemize}
        \item кнопка \texttt{Work with AI};
        \item форма настройки \texttt{.yml};
        \item виджет Recommendations на вкладке Code;
        \item страница detailed analysis;
        \item flow diagram для issue-to-PR;
        \item flow diagram для recommendations.
    \end{itemize}

    \item Поднять frontend skeleton:
    \begin{itemize}
        \item базовая структура frontend;
        \item routing;
        \item demo page, имитирующая repository/code page GitFlame;
        \item компонент \texttt{Work with AI};
        \item компонент recommendations widget;
        \item страница detailed recommendations.
    \end{itemize}

    \item Подключить demo frontend к mock/backend API:
    \begin{itemize}
        \item запрос статуса анализа;
        \item запрос summary;
        \item запрос списка recommendations;
        \item close/delete recommendation.
    \end{itemize}

    \item Реализовать recommendation endpoints:
    \begin{itemize}
        \item \texttt{GET /repositories/\{id\}/recommendations/status};
        \item \texttt{GET /repositories/\{id\}/recommendations/summary};
        \item \texttt{GET /repositories/\{id\}/recommendations};
        \item \texttt{PATCH /recommendations/\{id\}/close};
        \item \texttt{DELETE /recommendations/\{id\}}.
    \end{itemize}

    \item Подготовить screenshots/GIFs:
    \begin{itemize}
        \item работающий frontend skeleton;
        \item recommendations widget;
        \item detailed analysis page;
        \item Swagger/API docs.
    \end{itemize}
\end{enumerate}

\textbf{Deliverables:}
\begin{itemize}
    \item report draft
    \item user stories
    \item acceptance criteria
    \item wireframes/mockups
    \item frontend skeleton
    \item recommendation UI
    \item recommendation endpoints
    \item screenshots/GIFs
\end{itemize}

\section{Проверка покрытия требований Week 2}

\begin{longtable}{|p{0.58\textwidth}|p{0.32\textwidth}|}
\hline
\textbf{Требование Week 2} & \textbf{Ответственный} \\
\hline
Course familiarization & Данил \\
\hline
Team \& roles documentation & Рома \\
\hline
Project/track statement & Рома \\
\hline
Problem statement & Рома \\
\hline
Industrial justification & Данил \\
\hline
Week-by-week roadmap & Рома \\
\hline
Milestones, dependencies, risks & Данил \\
\hline
Kanban backlog & Данил \\
\hline
User stories & Рома \\
\hline
Acceptance criteria & Рома \\
\hline
MVP scope & Данил \\
\hline
Boilerplate structure & Артур \\
\hline
Runnable app & Артур \\
\hline
Docker support & Руслан \\
\hline
Tech stack justification & Данил \\
\hline
Wireframes/mockups & Рома \\
\hline
User flow diagrams & Рома \\
\hline
Frontend skeleton/routing & Рома \\
\hline
Backend API contract & Артур \\
\hline
Backend structure & Артур \\
\hline
Initial DB schema & Амир \\
\hline
ML experimentation setup & Карим \\
\hline
Core MVP journey & Артур \\
\hline
Frontend/backend connection & Рома \\
\hline
Data stored/retrieved & Амир \\
\hline
Open-source model selection for code generation & Данил \\
\hline
Open-source model selection for recommendations & Карим \\
\hline
Initial baseline/measurement & Карим \\
\hline
Internal demo \& feedback & Данил \\
\hline
Screenshots/GIFs & Рома \\
\hline
PR/Issue/API links & Рома \\
\hline
\end{longtable}

\section{Ожидаемый результат Sprint 1}

К концу Sprint 1 команда должна иметь:

\begin{itemize}
    \item структурированный репозиторий с backend, frontend, docs и Docker support;
    \item runnable backend endpoint;
    \item базовый frontend skeleton;
    \item OpenAPI/Swagger документацию;
    \item initial DB schema;
    \item mock ML service для автогенерации и рекомендаций;
    \item выбранные open-source модели для двух ML-сценариев;
    \item описанные prompts и JSON contracts;
    \item mock Git workflow для branch/PR;
    \item draft спецификацию \texttt{.yml};
    \item API-контракты для взаимодействия с внешним GitFlame;
    \item wireframes, user stories, acceptance criteria и MVP scope;
    \item Kanban backlog;
    \item материалы для отчёта, screenshots/GIFs и notes from internal demo.
\end{itemize}

\section{Минимальный end-to-end MVP flow}

В Sprint 1 необходимо показать хотя бы один рабочий end-to-end сценарий:

\begin{enumerate}
    \item Пользователь создаёт issue.
    \item GitFlame передаёт backend данные issue, \texttt{.yml} и repo context через API/webhook.
    \item Backend валидирует \texttt{.yml}.
    \item Backend отправляет issue, \texttt{.yml} и repo context в ML service.
    \item ML service возвращает \texttt{plan.md}.
    \item Пользователь подтверждает план.
    \item Backend формирует mock Git workflow response.
    \item Система возвращает GitFlame structured response с branch name, PR URL и reviewer.
\end{enumerate}

Дополнительный демонстрационный flow для рекомендаций:

\begin{enumerate}
    \item GitFlame передаёт системе найденный \texttt{.yml} и repo context.
    \item Запускается анализ репозитория.
    \item ML service возвращает summary и recommendation cards.
    \item Backend сохраняет результат в базе данных.
    \item Frontend отображает summary в виджете Recommendations.
    \item Пользователь открывает detailed analysis page.
\end{enumerate}

\end{document}
